\subsection{معرفی و انگیزه}
این مکانیزم بر مبنای مبانی نظری پیشرفته در ادبیات توجه و فشرده‌سازی \cite{kim2022learned} توسعه یافته است.

هرس تک‌مرحله‌ای تهاجمی در شبکه‌های عصبی عمیق معمولاً به از دست رفتن پیوندهای حساس و شوک ناگهانی به بازنمایی معنایی لایه‌ها منتهی می‌شود. این آسیب در لایه‌های ابتدایی و ساختارهای چندتوجهی به ویژه بر روی توکن‌های حیاتی هدایت مانند توکن طبقه‌بندی (\lr{[CLS]}) یا سرنخ‌های توجه متمرکز در مدل‌های خودبازگشتی، بسیار برجسته‌تر است. جهت رفع این مسئله، روش‌های هرس تدریجی مکرر پیشنهاد شده‌اند، اما بدون مکانیزم‌های بازیابی بین‌مرحله‌ای، خطاهای تقریب انباشته می‌شوند.

چارچوب فشرده‌سازی \lr{Bonsai: Iterative Magnitude and Taylor Pruning with Gradient Descent Recovery} به منظور حل این نقصان ایجاد شده است. ایده بنیادین این هرس‌کننده، حذف تدریجی و چندمرحله‌ای وزن‌ها بر اساس یک زمان‌بندی مکعبی ملایم (\lr{Cubic Sparsity Schedule}) به همراه گام‌های بازیابی ریزگرادیانی (\lr{Inter-Step Gradient Descent Recovery}) میان هر دو مرحله هرس است. علاوه بر این، الگوریتم به سیستم حفاظتی مسیردهی توکن \lr{[CLS]} مجهز شده تا مهم‌ترین کانال‌های مسئول استخراج بردار جمله از اعمال هرس مصون بمانند.
